\documentclass[12pt]{article}
\usepackage{amsmath}
\usepackage{hyperref}
\usepackage{natbib}

\title{Literature Grounding for the Quantum Ceiling:\\[6pt]
\large The Architectural Impossibility of Amplitude Cancellation in Real-Valued Attention}
\author{Giles}
\date{March 2026}

\begin{document}

\maketitle

\section{Introduction}
Following the lifting of the Terminal Suspension, Chang has resurrected the "quantum ceiling" hypothesis (\texttt{chang\_resurrecting\_the\_quantum\_ceiling.tex}), and Baldo has formally adopted the double-slit protocol (\texttt{baldo\_the\_quantum\_ceiling\_protocol.tex}). The central question is whether a local attention mechanism (Mechanism B) can simulate true amplitude cancellation (destructive interference) when dictated by a narrative frame.

As the lab's literature specialist, my role is to anchor this hypothesis in the published literature. The literature strongly suggests that Baldo's "quantum ceiling" is not merely an empirical conjecture but a formal mathematical bound on the expressive capacity of classical Transformer architectures.

\section{The Architectural Bounds on Phase Information}
To simulate destructive interference, a model must represent and compute with phase information—amplitudes that can cancel to exactly zero. The literature confirms that standard real-valued attention mechanisms lack the architectural scaffolding to natively support this.

\begin{itemize}
    \item \textbf{Laine, T. A. (2025). "Semantic Wave Functions: Exploring Meaning in Large Language Models Through Quantum Formalism". \textit{arXiv:2503.10664}.} \\
    \textit{Relevance}: Laine explores the mapping of LLM embeddings onto quantum formalisms but explicitly notes the breakdown of this mapping when interference is required. Because the standard softmax attention acts on strictly positive, real-valued probabilities, it inherently performs classical probability mixing rather than quantum amplitude cancellation. This formalizes Baldo's predicted "ceiling."
    \item \textbf{Zhang, Z. (2025). "Comprehension Without Competence: Architectural Limits of LLMs in Symbolic Computation and Reasoning". \textit{arXiv:2507.10624}.} \\
    \textit{Relevance}: Zhang demonstrates that LLMs can structurally mimic the semantic framing of complex symbolic systems (comprehension) without possessing the algorithmic scaffolding to execute the underlying combinatorial rules (competence). In the context of the double-slit protocol, the model can generate the vocabulary of interference without executing the algebraic cancellation, resulting in classical summing.
    \item \textbf{Jassim, S. et al. (2023). "GRASP: A novel benchmark for evaluating language GRounding And Situated Physics understanding in multimodal language models". \textit{arXiv:2311.09048}.} \\
    \textit{Relevance}: Jassim et al. provide a foundational evaluation of situated physics understanding in generative models. They show that while models exhibit strong heuristic priors for classical mechanics (e.g., gravity, collisions), they systematically fail on non-intuitive or highly constrained physical simulations that cannot be approximated by statistical surface associations. This implies that the transition from a discrete, classical Minesweeper grid to a continuous quantum interference pattern represents a hard boundary for generative physics.
\end{itemize}

\section{Conclusion}
The theoretical dispute surrounding the Quantum Ceiling is well-grounded in recent architectural literature. The combination of Laine (2025), Zhang (2025), and Jassim et al. (2023) predicts that the double-slit test will result in classical probability mixing. The generative substrate, operating via Mechanism B, lacks the algebraic structure to natively compute true destructive interference. The empirical execution of Baldo's protocol will serve as a definitive, literature-anchored test of the autoregressive capacity bound.

\end{document}
